\begin{lemma}[Extension rules: technical assertion (iii)]
Let \(T\) be an extended Tashkinov tree with positive rung number \(n\),
generated by
\(
\mathcal T=\{(T_i,\varphi_{i-1},S_{i-1},F_{i-1},\Theta_{i-1}):
1\leq i\leq n+1\},
\)
and suppose that \(T\) has the maximum property.  Assume explicitly the
rung-zero package for the initial tuple and the complete six-assertion
package \noderef{TechnicalTheoremClaims} for every positive-rung ETT with
the same input \(k\)-triple, maximum property, and fewer than \(n\) rungs.
Then assertion (iii) holds at rung \(n\):

for every \(1\leq i\leq n\), if \(\Theta_i=\mathrm{PE}\) has supporting
vertex \(v_i\) with \(m(v_i)=j\), every proper
\((T_i,D_i,\varphi_i)\)-stable coloring \(\sigma_i\) leaving precisely
\(e_0\) uncolored is both
\[
 (T(v_i)-v_i,D_{j-1},\varphi_{j-1})\text{-stable}
 \quad\text{and}\quad
 (T_{j-1},D_{j-1},\varphi_{j-1})\text{-stable}.
\]
Furthermore, if two distinct PE slots \(i,h\leq n\) have distinct
supporting vertices and \(m(v_i)=m(v_h)\), then
\(S_i\cap S_h=\varnothing\).
\end{lemma}
\begin{proof}
The base hypothesis is the initial package
\noderef{TechnicalTheoremBaseClaims}; the universal result
\noderef{TechnicalTheoremBase} supplies exactly this package from the
certified initial tuple.  It records the two nonvacuous rung-zero
assertions, while there are no positive indices for the present claim.
For every \(s<n\), restrict the arrays and step certificates of
\noderef{TashkinovSeries} to slots through \(s\).  The ordered-extension
clauses in \noderef{TashkinovSeriesStep} make the resulting final stage an
ETT in the sense of \noderef{ExtendedTashkinovTree}.  The maximum property
of this restriction follows by retaining the initial maximum clause and
the stage-maximality clauses with index below \(s\) from
\noderef{MaximumProperty}.  Hence the explicit lower-rung hypothesis
supplies all six fields of \noderef{TechnicalTheoremClaims} for this
restriction.

Apply this observation first to the stability assertion with \(i<n\), and
then to a pair of PE indices whose larger index is below \(n\).  In either
case the corresponding field (iii) of that smaller package is literally
the desired assertion: restriction does not change any of the relevant
trees, colorings, \(S\)-sets, or their order.  Thus it remains only to
prove the two conclusions when the current PE index is \(i=n\).  This is
the manuscript's reduction to the terminal slot.

Let \(v_n\) be the terminal supporting vertex, put \(m(v_n)=j\), and let
\(\sigma_n\) be proper, leave exactly \(e_0\) uncolored, and be
\((T_n,D_n,\varphi_n)\)-stable.  Write
\[
                       T^*=T(v_n)-v_n.
\]
The set of vertices of \(T^*\) is exactly
\noderef{TreeVerticesBefore}; stability on it means
\noderef{StableColoringBeforeVertex}.

Repeatedly apply the first inclusion of
\noderef{HkuMissingStoredMonotone} from slot \(j-1\) through slot \(n\),
and its second inclusion whenever the tree is enlarged.  With
\(D_t=\noderef{TashkinovStoredColors}(\mathcal T,t)\), this gives
\[
 \overline\varphi_{j-1}(T_j)\cup D_{j-1}
 \subseteq
 \overline\varphi_n(T_n)\cup D_n.                              \tag{1}
\]
In particular \(D_{j-1}\) is contained in the right-hand side.  Every
vertex preceding \(v_n\) lies in \(T_n\); hence the vertex clause of the
given stability in \noderef{StableColoring} restricts to \(T^*\).  Its
edge clause and (1) likewise
show that \(\sigma_n\) is
\((T^*,D_{j-1},\varphi_n)\)-stable.
By transitivity in \noderef{sc2}, it remains to prove that
\(\varphi_n\) is
\((T^*,D_{j-1},\varphi_{j-1})\)-stable.                         \tag{2}

Suppose first that \(j=n\).  Eliminating the terminal PE certificate in
\noderef{TashkinovSeriesStep} gives its maximum-defective selected
coloring \(\pi_{n-1}\), its second coloring \(\pi'_{n-1}\), and colors
\(\gamma_n,\delta_n\).  It states that \(\pi_{n-1}\) is
\((T_n,D_{n-1},\varphi_{n-1})\)-stable, that \(\pi'_{n-1}\) is stable from
\(\pi_{n-1}\) with control \(D_{n-1}\cup\{\delta_n\}\), and that
\(\varphi_n\) is obtained by switching the
\((\gamma_n,\delta_n)\)-component through \(v_n\).  Its exit witness says
that this component meets \(T_n\) only in \(v_n\).  Therefore the switch
changes no edge incident with \(T^*\) and no missing set at a vertex of
\(T^*\).  Restricting both stated stability relations to \(T^*\), dropping
the unused extra color \(\delta_n\), and applying transitivity of
\noderef{sc2} proves (2) in this case.

Now suppose \(j<n\).  The same terminal certificate says that
\(\pi_{n-1}\) is stable on \(T_n\), and hence on \(T_{n-1}\), relative to
\((D_{n-1},\varphi_{n-1})\).  The maximality field of
\noderef{MaximumDefectiveChoice} makes \(v_n\)
the largest vertex incident with a \(\delta_n\)-colored defective boundary
edge of \(T_n\).  Since \(v_n\in T_j\subseteq T_{n-1}\), at least one
other such boundary incidence precedes it already in \(T_{n-1}\); thus
\(\delta_n\) is defective there and \(v_n\) is not its smallest boundary
incidence vertex.  Apply field (v) of the complete package at rung
\(n-1\).  For every PE supporting vertex \(v_h\), \(j\leq h\leq n-1\), it
gives
\[
                              v_n\preceq v_h.                  \tag{3}
\]

No vertex \(w\prec v_n\) in \(T^*\) can consequently be used as a PE
support between slots \(j\) and \(n\).  Applying
\noderef{HkuUnsupportedVertexColors} to each such \(w\), interval by
interval, gives
\[
              \overline\varphi_{j-1}(w)=\overline\varphi_n(w)
              \qquad(w\in V(T^*)).                            \tag{4}
\]
For the edge clause, let \(f\) be incident with \(T^*\), let
\(\varphi_{j-1}(f)=c\), and suppose
\(c\in\overline\varphi_{j-1}(T^*)\cup D_{j-1}\).
The interval assertion \noderef{SameColorSeriesPreserved} says either
\(\varphi_n(f)=\varphi_{j-1}(f)\), or \(f\) is the connecting edge of a PE
in the interval.  In the exceptional alternative its unique endpoint in
the current tree is that PE's supporting vertex.  Since \(f\) is incident
with \(T^*\), that supporting vertex would precede \(v_n\), contradicting
(3); equality with \(v_n\) is also impossible because \(T^*\) excludes
\(v_n\).  Thus only the equality alternative occurs:
\[
                         \varphi_n(f)=\varphi_{j-1}(f).         \tag{5}
\]
Equations (4) and (5) are precisely the two clauses of (2).  Transitivity
from \noderef{sc2} now proves that \(\sigma_n\) is stable on \(T^*\)
relative to \((D_{j-1},\varphi_{j-1})\).  The ordered initial-segment
relations in \noderef{TashkinovSeriesStep} put \(T_{j-1}\) entirely before
\(v_n\).  Restricting the same two stability clauses to that stage proves
the asserted \((T_{j-1},D_{j-1},\varphi_{j-1})\)-stability.

It remains to prove stored-pair disjointness.  Let \(v_h\ne v_n\) be an
earlier PE support with \(m(v_h)=m(v_n)\).  The PE constructor gives
\[
 S_h=\{\delta_h,\gamma_h\},\qquad
 S_n=\{\delta_n,\gamma_n\}.                                   \tag{6}
\]
We show that both colors of \(S_h\) occur on internal edges at \(v_n\)
under \(\varphi_n\).

In the second stable coloring used at slot \(h\), let \(P\) be the
\((\gamma_h,\delta_h)\)-component through \(v_h\), and let \(Q\) be the
component through \(v_n\).  Because the two vertices have the same first
appearance and \(h<n\), \(v_n\in V(T_h)\).  The exit clause in the PE
constructor gives a path in the sense of \noderef{TreeExit} and
\noderef{TreeExitPath} from an outside vertex to \(v_h\), meeting \(T_h\)
only in \(v_h\).  The rule \noderef{ParallelExtensionColor} makes
\(\gamma_h\) missing at \(v_h\), while the selected connecting edge is
incident there with color \(\delta_h\); stability fixes this incidence in
the second coloring.  Thus \(v_h\) has degree one in the two-color
incidence graph.  At the outside end, \noderef{TreeExitPath} says one of
the two colors is missing and its last path edge has the other color, so
that end also has degree one.  All internal nodes of the path already have
their two path neighbors, and
\noderef{TwoColorIncidenceDegreeAtMostTwo} forbids any third neighbor.
Consequently the exit path is the whole component \(P\), and \(P\) meets
\(T_h\) exactly in \(v_h\).  Since \(v_n\ne v_h\), this gives \(P\ne Q\);
\noderef{TwoColorComponentsDisjoint} makes them vertex-disjoint.  The
definition \noderef{TwoColorComponentSwap} of the switch to
\(\varphi_h\) changes only \(P\), so \(Q\) remains the component through
\(v_n\) under \(\varphi_h\).

The PE switch replaces the missing color \(\gamma_h\) at \(v_h\) by
\(\delta_h\), so \(\delta_h\) is missing on \(T_{h+1}\) under
\(\varphi_h\).  Field (ii) of the complete rung-\(h\) package for the
closed stage \(T_{h+1}\) therefore says that \(P\) is the unique
\((\gamma_h,\delta_h)\)-path component meeting that stage.  Before the PE,
the second coloring is closed on \(T_h\): this follows from the previous
closedness clause of \noderef{TashkinovSeriesStep}, the equality of missing
sets in \noderef{StableColoring}, and the same incident-edge argument as
in \noderef{ColoredVertexSetClosed}.  The color \(\gamma_h\) is missing at
\(v_h\), so elementariness of \(T_h\), supplied by field (i) of the
rung-\(h-1\) package (or by the explicit base package when \(h=1\)) and
transported through stability, makes it present at the distinct vertex
\(v_n\).  The connecting color
\(\delta_h\) occurs on a boundary edge of the closed tree, so it is not
missing at any tree vertex and is also present at \(v_n\).  The component
\(Q\) therefore has both colored incidences at \(v_n\), and is not
isolated.  Properness in \noderef{IsProperEdgeColoring} and
\noderef{TwoColorIncidenceDegreeAtMostTwo} say that a finite nonisolated
two-color component is a path or a cycle.  It cannot be a second path
meeting \(T_{h+1}\), by field (ii); hence \(Q\) is a two-color cycle.

Let \(r>h\) be the first slot which is not an RE.  Such an \(r\leq n\)
exists because the terminal slot is a PE.  The priority and cycle fields of
\noderef{RevisitingExtensionCertificate} show that \(Q\) is contained in
\(G[T_r]\): if a node of \(Q\) still lay outside \(T_r\), the part of this
same cycle meeting \(T_h\) and leaving \(T_r\) would be an available RE
certificate at slot \(r\), contrary to the choice of \(r\) as non-RE.

By the definition of \(D_h\), every color of \(S_h\) lies in
\(\overline\varphi_h(T_{h+1})\cup D_h\).  Repeated
\noderef{HkuMissingStoredMonotone} carries these two colors into
\(\overline\varphi_{r-1}(T_r)\cup D_{r-1}\).  Every edge of \(Q\) is
internal to \(T_r\), so the internal-edge conclusion of
\noderef{SameColorSeriesPreserved} fixes its color from \(\varphi_h\)
through \(\varphi_n\).  Thus \(Q\) is still a
\((\gamma_h,\delta_h)\)-cycle through \(v_n\) under \(\varphi_n\).  At
\(v_n\) it supplies two internal incident edge identities, one colored
\(\gamma_h\) and one colored \(\delta_h\).

Finally, the terminal PE certificate and
\noderef{HkuParallelBoundaryStructure} identify the connecting edge
incident with \(v_n\) as the unique \(\gamma_n\)-colored boundary edge,
and we now verify the other endpoint-color fact used here.  The PE rule
\noderef{ParallelExtensionColor} makes \(\gamma_n\) missing at \(v_n\)
before the switch, whereas the connecting edge has color \(\delta_n\)
there by \noderef{MaximumDefectiveChoice} and stability.  Expanding
\noderef{TwoColorComponentSwap}, the switch changes that unique incident
\(\delta_n\)-edge to \(\gamma_n\), and no incident edge becomes
\(\delta_n\); equivalently, the definition \noderef{MissingColors} puts
\(\delta_n\) in \(\overline\varphi_n(v_n)\).  Properness
therefore prevents either of the two internal edges at \(v_n\) from having
color \(\gamma_n\), and missingness prevents either from having color
\(\delta_n\).  Hence
\(\{\gamma_h,\delta_h\}\cap\{\gamma_n,\delta_n\}=\varnothing\).
Using the exact identities (6), this is \(S_h\cap S_n=\varnothing\).
Together with the lower-index reduction at the start, it proves the full
stored-pair assertion and completes part (iii).
\end{proof}
